\section{第 9 章综合习题}

\begin{problem}{多元函数微分学}{Partial-Derivative-Condition}
    设 $a_1, a_2, \cdots, a_n$ 是非零常数。$f(x_1, x_2, \cdots, x_n)$ 在 $\mathbb{R}^n$ 上可微。求证：存在 $\mathbb{R}$ 上一元可微函数 $F(s)$ 使得 $f(x_1, x_2, \cdots, x_n) = F(a_1 x_1 + a_2 x_2 + \cdots + a_n x_n)$ 的充分必要条件是 $a_j \frac{\partial f}{\partial x_i} = a_i \frac{\partial f}{\partial x_j}, i, j = 1, 2, \cdots, n$。
\end{problem}

\begin{solution}

    \textbf{必要性}：
    设 $s = \sum_{k=1}^n a_k x_k$，且 $f(x_1, \dots, x_n) = F(s)$。
    由链式法则，对任意 $i \in \{1, \dots, n\}$，有
    \begin{equation}
        \frac{\partial f}{\partial x_i} = F'(s) \frac{\partial s}{\partial x_i} = a_i F'(s).
    \end{equation}
    同理，对任意 $j \in \{1, \dots, n\}$，有
    \begin{equation}
        \frac{\partial f}{\partial x_j} = a_j F'(s).
    \end{equation}
    从而
    \begin{equation}
        a_j \frac{\partial f}{\partial x_i} = a_j a_i F'(s) = a_i a_j F'(s) = a_i \frac{\partial f}{\partial x_j}.
    \end{equation}
    必要性得证。

    \textbf{充分性}：
    已知 $a_j \frac{\partial f}{\partial x_i} = a_i \frac{\partial f}{\partial x_j}$ 对所有 $i, j$ 成立。
    由于 $a_k \neq 0$，可变形为
    \begin{equation}
        \frac{1}{a_i} \frac{\partial f}{\partial x_i} = \frac{1}{a_j} \frac{\partial f}{\partial x_j}.
    \end{equation}
    这表明 $\frac{1}{a_k} \frac{\partial f}{\partial x_k}$ 的值与下标 $k$ 无关。记该公共值为 $g(x_1, \dots, x_n)$，即
    \begin{equation}
        \frac{\partial f}{\partial x_k} = a_k g(x_1, \dots, x_n), \quad k=1, \dots, n.
    \end{equation}
    考虑 $f$ 的全微分：
    \begin{equation}
        \d f = \sum_{k=1}^n \frac{\partial f}{\partial x_k} \d x_k = \sum_{k=1}^n a_k g \d x_k = g \sum_{k=1}^n a_k \d x_k = g \d \left( \sum_{k=1}^n a_k x_k \right).
    \end{equation}
    令 $s = \sum_{k=1}^n a_k x_k$。上式表明 $\d f = g \d s$。
    这意味着 $f$ 的全微分仅依赖于 $\d s$，即 $f$ 的变化完全由 $s$ 的变化决定。
    在单连通区域 $\mathbb{R}^n$ 上，这意味着 $f$ 仅是 $s$ 的函数。
    即存在一元可微函数 $F$，使得
    \begin{equation}
        f(x_1, \dots, x_n) = F(s) = F(a_1 x_1 + a_2 x_2 + \cdots + a_n x_n).
    \end{equation}
    充分性得证。
\end{solution}

v

\begin{problem}{多元复合函数求导}{Partial-Derivative-Chain-Rule-3}
    设 $f(x,y,z) = F(u,v,w)$，其中 $x^2 = vw$、$y^2 = wu$、$z^2 = uv$。求证：
    \begin{equation}
        x \frac{\partial f}{\partial x} + y \frac{\partial f}{\partial y} + z \frac{\partial f}{\partial z} = u \frac{\partial F}{\partial u} + v \frac{\partial F}{\partial v} + w \frac{\partial F}{\partial w}.
    \end{equation}
\end{problem}

\begin{myproof}
    将 $x^2 = vw$、$y^2 = wu$、$z^2 = uv$ 关于 $x$ 求偏导，得
    \begin{equation}
        \begin{cases}
            2x = v \frac{\partial w}{\partial x} + w \frac{\partial v}{\partial x} \\
            0 = w \frac{\partial u}{\partial x} + u \frac{\partial w}{\partial x}  \\
            0 = u \frac{\partial v}{\partial x} + v \frac{\partial u}{\partial x}
        \end{cases}
        \label{eq:differentiation-system}
    \end{equation}
    由 \cref{eq:differentiation-system} 后两式可得 $\frac{\partial w}{\partial x} = -\frac{w}{u} \frac{\partial u}{\partial x}$ 且 $\frac{\partial v}{\partial x} = -\frac{v}{u} \frac{\partial u}{\partial x}$。代入第一式，得
    \begin{equation}
        2x = v \left( -\frac{w}{u} \frac{\partial u}{\partial x} \right) + w \left( -\frac{v}{u} \frac{\partial u}{\partial x} \right) = -\frac{2vw}{u} \frac{\partial u}{\partial x} = -\frac{2x^2}{u} \frac{\partial u}{\partial x}.
    \end{equation}
    由此解得 $\frac{\partial u}{\partial x} = -\frac{u}{x}$，进而可得 $\frac{\partial v}{\partial x} = \frac{v}{x}$、$\frac{\partial w}{\partial x} = \frac{w}{x}$。由对称性可得所有偏导项：
    \begin{equation}
        \begin{aligned}
            \frac{\partial u}{\partial x} = -\frac{u}{x}, \quad \frac{\partial v}{\partial x} = \frac{v}{x}, \quad \frac{\partial w}{\partial x} = \frac{w}{x}; \\
            \frac{\partial u}{\partial y} = \frac{u}{y}, \quad \frac{\partial v}{\partial y} = -\frac{v}{y}, \quad \frac{\partial w}{\partial y} = \frac{w}{y}; \\
            \frac{\partial u}{\partial z} = \frac{u}{z}, \quad \frac{\partial v}{\partial z} = \frac{v}{z}, \quad \frac{\partial w}{\partial z} = -\frac{w}{z}.
        \end{aligned}
    \end{equation}
    根据复合函数求导法则，有
    \begin{equation}
        \begin{aligned}
            x \frac{\partial f}{\partial x} & = x \left( \frac{\partial F}{\partial u} \frac{\partial u}{\partial x} + \frac{\partial F}{\partial v} \frac{\partial v}{\partial x} + \frac{\partial F}{\partial w} \frac{\partial w}{\partial x} \right) = -u \frac{\partial F}{\partial u} + v \frac{\partial F}{\partial v} + w \frac{\partial F}{\partial w}, \\
            y \frac{\partial f}{\partial y} & = y \left( \frac{\partial F}{\partial u} \frac{\partial u}{\partial y} + \frac{\partial F}{\partial v} \frac{\partial v}{\partial y} + \frac{\partial F}{\partial w} \frac{\partial w}{\partial y} \right) = u \frac{\partial F}{\partial u} - v \frac{\partial F}{\partial v} + w \frac{\partial F}{\partial w},  \\
            z \frac{\partial f}{\partial z} & = z \left( \frac{\partial F}{\partial u} \frac{\partial u}{\partial z} + \frac{\partial F}{\partial v} \frac{\partial v}{\partial z} + \frac{\partial F}{\partial w} \frac{\partial w}{\partial z} \right) = u \frac{\partial F}{\partial u} + v \frac{\partial F}{\partial v} - w \frac{\partial F}{\partial w}.
        \end{aligned}
    \end{equation}
    将上述三式相加，整理可得
    \begin{equation}
        \boxed{x \frac{\partial f}{\partial x} + y \frac{\partial f}{\partial y} + z \frac{\partial f}{\partial z} = u \frac{\partial F}{\partial u} + v \frac{\partial F}{\partial v} + w \frac{\partial F}{\partial w}}
    \end{equation}
\end{myproof}


\begin{problem}{齐次函数欧拉定理}{Euler-Theorem-Homogeneous}
    若函数 $u = f(x, y, z)$ 满足恒等式 $f(tx, ty, tz) = t^k f(x, y, z)$ ($t > 0$)，则称 $f(x, y, z)$ 为 $k$ 次齐次函数。试证下述关于齐次函数的欧拉定理：可微函数 $f(x, y, z)$ 为 $k$ 次齐次函数的充要条件是：
    \begin{equation}
        x f'_x(x, y, z) + y f'_y(x, y, z) + z f'_z(x, y, z) = k f(x, y, z).
    \end{equation}
\end{problem}

\begin{myproof}
    必要性：若 $f(x, y, z)$ 是 $k$ 次齐次函数，则有恒等式 $f(tx, ty, tz) = t^k f(x, y, z)$。对此式两端关于 $t$ 求导，由多元复合函数求导法则得：
    \begin{equation}
        \begin{aligned}
            \frac{\d}{\d t} f(tx, ty, tz) & = \frac{\partial f}{\partial (tx)} \cdot \frac{\d (tx)}{\d t} + \frac{\partial f}{\partial (ty)} \cdot \frac{\d (ty)}{\d t} + \frac{\partial f}{\partial (tz)} \cdot \frac{\d (tz)}{\d t} \\
                                          & = x f'_x(tx, ty, tz) + y f'_y(tx, ty, tz) + z f'_z(tx, ty, tz).
        \end{aligned} \label{eq:Diff-LHS}
    \end{equation}
    同时，等式右端关于 $t$ 的导数为：
    \begin{equation}
        \frac{\d}{\d t} \left[ t^k f(x, y, z) \right] = k t^{k-1} f(x, y, z). \label{eq:Diff-RHS}
    \end{equation}
    联立 \cref{eq:Diff-LHS} 与 \cref{eq:Diff-RHS}，并令 $t = 1$，即得：
    \begin{equation}
        x f'_x(x, y, z) + y f'_y(x, y, z) + z f'_z(x, y, z) = k f(x, y, z).
    \end{equation}

    充分性：若已知题设微分方程成立。对于任意固定的 $(x, y, z)$，构造辅助函数：
    \begin{equation}
        G(t) = \frac{f(tx, ty, tz)}{t^k} = t^{-k} f(tx, ty, tz) \quad (t > 0).
    \end{equation}
    对 $G(t)$ 关于 $t$ 求导得：
    \begin{equation}
        \begin{aligned}
            G'(t) & = -k t^{-k-1} f(tx, ty, tz) + t^{-k} \frac{\d}{\d t} f(tx, ty, tz)                                                           \\
                  & = -k t^{-k-1} f(tx, ty, tz) + t^{-k} \left[ x f'_x(tx, ty, tz) + y f'_y(tx, ty, tz) + z f'_z(tx, ty, tz) \right]             \\
                  & = -k t^{-k-1} f(tx, ty, tz) + t^{-k-1} \left[ (tx) f'_x(tx, ty, tz) + (ty) f'_y(tx, ty, tz) + (tz) f'_z(tx, ty, tz) \right].
        \end{aligned}
    \end{equation}
    由于题设等式对任意点均成立，将点 $(x, y, z)$ 替换为 $(tx, ty, tz)$ 可得：
    \begin{equation}
        (tx) f'_x(tx, ty, tz) + (ty) f'_y(tx, ty, tz) + (tz) f'_z(tx, ty, tz) = k f(tx, ty, tz).
    \end{equation}
    代入 $G'(t)$ 的表达式中，有：
    \begin{equation}
        G'(t) = -k t^{-k-1} f(tx, ty, tz) + t^{-k-1} \left[ k f(tx, ty, tz) \right] = 0.
    \end{equation}
    因此 $G(t)$ 为常数。由 $G(1) = f(x, y, z)$ 可知 $G(t) = f(x, y, z)$，即：
    \begin{equation}
        \frac{f(tx, ty, tz)}{t^k} = f(x, y, z) \implies f(tx, ty, tz) = t^k f(x, y, z).
    \end{equation}
    综上所述，原命题得证。
    \begin{equation}
        \boxed{x f'_x(x, y, z) + y f'_y(x, y, z) + z f'_z(x, y, z) = k f(x, y, z)}
    \end{equation}
\end{myproof}


\begin{problem}{齐次函数与隐函数}{Homogeneous-Implicit-Function}
    设 $f(x,y,z)$ 是 $n$ 次齐次的可微函数。若方程 $f(x,y,z) = 0$ 隐含函数 $z = \varphi(x,y)$ (即，$f'_z \neq 0$), 则 $\varphi(x,y)$ 是一次齐次函数。
\end{problem}

\begin{myproof}
    由于 $f(x, y, z)$ 是 $n$ 次齐次函数，根据齐次函数的定义，对于任意 $t > 0$，恒有
    \begin{equation}
        f(tx, ty, tz) = t^n f(x, y, z). \label{eq:homo-def}
    \end{equation}
    由题意，$z = \varphi(x, y)$ 是方程 $f(x, y, z) = 0$ 确定的隐函数，即满足
    \begin{equation}
        f(x, y, \varphi(x, y)) \equiv 0. \label{eq:implicit-relation}
    \end{equation}
    在 \cref{eq:homo-def} 中，令 $z = \varphi(x, y)$，并结合 \cref{eq:implicit-relation} 可得
    \begin{equation}
        f(tx, ty, t\varphi(x, y)) = t^n f(x, y, \varphi(x, y)) = t^n \cdot 0 = 0. \label{eq:scaled-zero}
    \end{equation}
    另一方面，对于点 $(tx, ty)$，隐函数 $z = \varphi(tx, ty)$ 同样满足方程
    \begin{equation}
        f(tx, ty, \varphi(tx, ty)) = 0. \label{eq:scaled-implicit}
    \end{equation}
    由于 $f'_z \neq 0$，根据隐函数存在定理，方程 $f(tx, ty, z) = 0$ 在该点附近确定的解 $z$ 是唯一的。通过比较 \cref{eq:scaled-zero} 与 \cref{eq:scaled-implicit} 可知
    \begin{equation}
        \varphi(tx, ty) = t\varphi(x, y).
    \end{equation}
    由此证明 $\varphi(x, y)$ 是一次齐次函数。
    \begin{equation}
        \boxed{\varphi(tx, ty) = t\varphi(x, y)}
    \end{equation}
\end{myproof}



\begin{problem}{多元函数方程}{Multivariate-Functional-Equation}
    设 $f(x,y)$ 在 $\mathbb{R}^2$ 上有连续二阶偏导数，且对任意实数 $x,y,z$ 满足 $f(x,y)=f(y,x)$ 和
    \begin{equation}
        f(x,y) + f(y,z) + f(z,x) = 3f\left(\frac{x+y+z}{3}, \frac{x+y+z}{3}\right).
    \end{equation}
    试求 $f(x,y)$。
\end{problem}

\begin{solution}
    \ansitem{1}

    记 $g(t) = f(t,t)$。对原方程两端关于 $x$ 求偏导得：
    \begin{equation}
        \frac{\partial f}{\partial x}(x,y) + \frac{\partial f}{\partial x}(z,x) = g'\left( \frac{x+y+z}{3} \right).
    \end{equation}
    由 $f(x,y) = f(y,x)$ 的对称性可知 $\frac{\partial f}{\partial x}(z,x) = \frac{\partial f}{\partial x}(x,z)$，代入上式得：
    \begin{equation}
        \frac{\partial f}{\partial x}(x,y) + \frac{\partial f}{\partial x}(x,z) = g'\left( \frac{x+y+z}{3} \right). \label{eq:partial-x-1}
    \end{equation}
    将 \cref{eq:partial-x-1} 两端关于 $y$ 求偏导：
    \begin{equation}
        \frac{\partial^2 f}{\partial x \partial y}(x,y) = \frac{1}{3} g''\left( \frac{x+y+z}{3} \right). \label{eq:second-partial-1}
    \end{equation}
    由于 \cref{eq:second-partial-1} 左端仅与 $x,y$ 相关，而右端包含变量 $z$，取 $z = -x-y$，$x$ 和 $y$ 自由变化时右侧为定值，故 $\frac{\partial^2 f}{\partial x \partial y}$ 必为常数。设该常数为 $4a$，则：
    \begin{equation}
        f(x,y) = 4axy + \phi(x) + \psi(y).
    \end{equation}
    由 $f(x,y) = f(y,x)$ 可知 $\phi(x) + \psi(y) = \phi(y) + \psi(x)$，由此推得 $\phi(x) - \psi(x) = C$。令 $h(x) = \phi(x) - \frac{C}{2}$，则 $f(x,y)$ 可表示为：
    \begin{equation}
        f(x,y) = 4axy + h(x) + h(y).
    \end{equation}
    此时 $g(t) = f(t,t) = 4at^2 + 2h(t)$。代入 \cref{eq:second-partial-1} 可得 $4a = \frac{1}{3}(8a + 2h''(t))$，解得 $h''(t) = 2a$。故 $h(x) = ax^2 + bx + c$。综上所述：
    \begin{equation}
        f(x,y) = 4axy + ax^2 + bx + ay^2 + by + C = a(x^2 + y^2 + 4xy) + b(x+y) + C.
    \end{equation}
    代入原方程验证，左端为：
    \begin{equation}
        \begin{aligned}
            \text{LHS} & = a(x^2+y^2+4xy + y^2+z^2+4yz + z^2+x^2+4zx) + 2b(x+y+z) + 3C \\
                       & = 2a(x^2+y^2+z^2 + 2xy + 2yz + 2zx) + 2b(x+y+z) + 3C          \\
                       & = 2a(x+y+z)^2 + 2b(x+y+z) + 3C.
        \end{aligned}
    \end{equation}
    右端为：
    \begin{equation}
        \begin{aligned}
            \text{RHS} & = 3 \left[ 6a \left( \frac{x+y+z}{3} \right)^2 + 2b \left( \frac{x+y+z}{3} \right) + C \right] \\
                       & = 2a(x+y+z)^2 + 2b(x+y+z) + 3C.
        \end{aligned}
    \end{equation}
    两端相等，故结果正确。
    \begin{equation}
        \boxed{f(x,y) = a(x^2 + y^2 + 4xy) + b(x + y) + C}
    \end{equation}
\end{solution}


\begin{problem}{多元函数不等式}{Multivariate-Inequality}
    证明：不等式 $\frac{x^2+y^2}{4} \leqslant \e^{x+y-2}$ ($x \geqslant 0, y \geqslant 0$)。
\end{problem}

\begin{myproof}
    构造辅助函数
    \begin{equation}
        g(x, y) = (x^2 + y^2) \e^{2-x-y}, \quad x \geqslant 0, y \geqslant 0.
    \end{equation}
    对 $g(x, y)$ 求偏导数，可得
    \begin{equation}
        \begin{aligned}
            \frac{\partial g}{\partial x} & = (2x - x^2 - y^2) \e^{2-x-y}, \\
            \frac{\partial g}{\partial y} & = (2y - x^2 - y^2) \e^{2-x-y}.
        \end{aligned}
    \end{equation}
    令偏导数均为零，由于指数项恒大于零，在第一象限内部需满足 $x^2 + y^2 = 2x$ 且 $x^2 + y^2 = 2y$。由此可解得区域内的唯一驻点为 $(1, 1)$。此时函数值为
    \begin{equation}
        g(1, 1) = 2 \e^{2-1-1} = 2.
    \end{equation}
    考虑区域边界。当 $x = 0$ 时，$g(0, y) = y^2 \e^{2-y}$。令其关于 $y$ 的导数 $(2y - y^2) \e^{2-y} = 0$，得临界点 $y = 0$ 或 $y = 2$。计算边界上的极值点：
    \begin{equation}
        g(0, 0) = 0, \quad g(0, 2) = 4.
    \end{equation}
    对称地，当 $y = 0$ 时，最大值为 $g(2, 0) = 4$。由于当 $x \to +\infty$ 或 $y \to +\infty$ 时，指数衰减速度远快于多项式增长，故 $g(x, y) \to 0$。综上可知，$g(x, y)$ 在闭区域 $x \geqslant 0, y \geqslant 0$ 上的全局最大值为 $4$。即满足
    \begin{equation}
        (x^2 + y^2) \e^{2-x-y} \leqslant 4.
    \end{equation}
    将不等式两边同除以 $4\e^{2-x-y}$，即得
    \begin{equation}
        \boxed{\frac{x^2+y^2}{4} \leqslant \e^{x+y-2}}.
    \end{equation}
\end{myproof}

\begin{problem}{多元函数连续性}{Multivariate-Continuity-2}
    设在 $\symbb{R}^3$ 上定义的 $u = f(x,y,z)$ 是 $z$ 的连续函数，且 $\frac{\partial f}{\partial x}$、$\frac{\partial f}{\partial y}$ 在 $\symbb{R}^3$ 上连续。证明：$u$ 在 $\symbb{R}^3$ 上连续。
\end{problem}

\begin{myproof}
    任取点 $P_0(x_0, y_0, z_0) \in \symbb{R}^3$。考察函数在该点的改变量
    \begin{equation}
        \Delta f = f(x, y, z) - f(x_0, y_0, z_0).
    \end{equation}
    将其分解为
    \begin{equation}
        \Delta f = [f(x, y, z) - f(x_0, y, z)] + [f(x_0, y, z) - f(x_0, y_0, z)] + [f(x_0, y_0, z) - f(x_0, y_0, z_0)].
    \end{equation}
    利用拉格朗日中值定理，对前两项分别有
    \begin{equation}
        \begin{aligned}
            f(x, y, z) - f(x_0, y, z)     & = \frac{\partial f}{\partial x}(\xi, y, z) (x - x_0), \quad \xi \text{ 介于 } x, x_0 \text{ 之间}；     \\
            f(x_0, y, z) - f(x_0, y_0, z) & = \frac{\partial f}{\partial y}(x_0, \eta, z) (y - y_0), \quad \eta \text{ 介于 } y, y_0 \text{ 之间}。
        \end{aligned}
    \end{equation}
    由于偏导数 $\frac{\partial f}{\partial x}$ 与 $\frac{\partial f}{\partial y}$ 在 $\symbb{R}^3$ 上连续，故在 $P_0$ 的任一紧邻域内必有界。设在该邻域内满足 $\left| \frac{\partial f}{\partial x} \right| \leqslant M$ 且 $\left| \frac{\partial f}{\partial y} \right| \leqslant M$，则有
    \begin{equation}
        |\Delta f| \leqslant M|x - x_0| + M|y - y_0| + |f(x_0, y_0, z) - f(x_0, y_0, z_0)|.
    \end{equation}
    当 $(x, y, z) \to (x_0, y_0, z_0)$ 时，显然有 $|x - x_0| \to 0$ 且 $|y - y_0| \to 0$。此外，由 $f$ 是关于 $z$ 的连续函数知
    \begin{equation}
        \lim_{z \to z_0} [f(x_0, y_0, z) - f(x_0, y_0, z_0)] = 0.
    \end{equation}
    根据夹逼定理，必有 $\lim_{(x, y, z) \to (x_0, y_0, z_0)} \Delta f = 0$，即
    \begin{equation}
        \boxed{\lim_{(x, y, z) \to (x_0, y_0, z_0)} f(x, y, z) = f(x_0, y_0, z_0)}.
    \end{equation}
    由于 $P_0$ 具有任意性，证得 $u = f(x, y, z)$ 在 $\symbb{R}^3$ 上连续。
\end{myproof}


\begin{problem}{常值函数判定}{Constant-Function-Convex-Domain}
    设 $D \subset \mathbb{R}^2$ 是包含原点的凸区域，$f \in C^1(D)$。证明：若
    \begin{equation}
        x \frac{\partial f(x,y)}{\partial x} + y \frac{\partial f(x,y)}{\partial y} = 0, \quad ((x,y) \in D), \label{eq:Euler-PDE}
    \end{equation}
    则 $f(x,y)$ 是常数。
\end{problem}

\begin{myproof}
    对任意 $(x, y) \in D$，构造辅助函数 $F(t) = f(tx, ty)$，$t \in [0, 1]$。由于 $D$ 是包含原点的凸区域，对于所有 $t \in [0, 1]$，点 $(tx, ty)$ 均属于 $D$。由复合函数求导法则有
    \begin{equation}
        F'(t) = x \frac{\partial f}{\partial x}(tx, ty) + y \frac{\partial f}{\partial y}(tx, ty).
    \end{equation}
    当 $t \in (0, 1]$ 时，由于 $(tx, ty) \in D$，利用 \cref{eq:Euler-PDE} 可得
    \begin{equation}
        F'(t) = \frac{1}{t} \left[ (tx) \frac{\partial f}{\partial x}(tx, ty) + (ty) \frac{\partial f}{\partial y}(tx, ty) \right] = 0.
    \end{equation}
    由于 $f \in C^1(D)$，其偏导数连续，故 $F'(t)$ 在 $[0, 1]$ 上连续。由 $F'(t) = 0$ 在 $(0, 1]$ 上成立及连续性可知，$F'(t) \equiv 0$ 对所有 $t \in [0, 1]$ 成立。因此
    \begin{equation}
        f(x, y) = F(1) = F(0) = f(0, 0).
    \end{equation}
    这表明 $f(x, y)$ 在 $D$ 上恒为常数。
    \begin{equation}
        \boxed{f(x, y) = C}
    \end{equation}
\end{myproof}


\begin{problem}{Hadamard 引理}{Hadamard-Lemma-R2}
    设 $f \in C^1(\mathbb{R}^2)$，$f(0,0) = 0$。证明：存在 $\mathbb{R}^2$ 上的连续函数 $g_1, g_2$ 使得
    \begin{equation}
        f(x,y) = x g_1(x,y) + y g_2(x,y).
    \end{equation}
\end{problem}

\begin{myproof}
    根据微积分基本定理，对于平面上任意一点 $(x, y)$，有
    \begin{equation}
        \begin{aligned}
            f(x, y) & = f(x, y) - f(0, 0)                        \\
                    & = \int_0^1 \frac{\d}{\d t} f(tx, ty) \d t.
        \end{aligned} \label{eq:FTC-Expansion}
    \end{equation}
    应用复合函数求导法则，被积函数可表示为
    \begin{equation}
        \frac{\d}{\d t} f(tx, ty) = x \frac{\partial f}{\partial x}(tx, ty) + y \frac{\partial f}{\partial y}(tx, ty).
    \end{equation}
    代入 \cref{eq:FTC-Expansion}，并利用积分的线性性质得
    \begin{equation}
        f(x, y) = x \int_0^1 \frac{\partial f}{\partial x}(tx, ty) \d t + y \int_0^1 \frac{\partial f}{\partial y}(tx, ty) \d t.
    \end{equation}
    令
    \begin{equation}
        g_1(x, y) = \int_0^1 \frac{\partial f}{\partial x}(tx, ty) \d t, \quad g_2(x, y) = \int_0^1 \frac{\partial f}{\partial y}(tx, ty) \d t.
    \end{equation}
    由于 $f \in C^1(\mathbb{R}^2)$，其偏导数 $\frac{\partial f}{\partial x}$ 与 $\frac{\partial f}{\partial y}$ 在 $\mathbb{R}^2$ 上连续。根据含参变量积分的连续性定理可知，$g_1(x, y)$ 与 $g_2(x, y)$ 在 $\mathbb{R}^2$ 上均为连续函数。由此得证
    \begin{equation}
        \boxed{f(x, y) = x g_1(x, y) + y g_2(x, y)}
    \end{equation}
\end{myproof}


\begin{problem}{全微分存在性充分条件}{Differentiability-Condition}
    设 $f(x, y)$ 在 $(x_0, y_0)$ 的某个邻域 $U$ 上有定义，$\frac{\partial f}{\partial x}$ 和 $\frac{\partial f}{\partial y}$ 在 $U$ 上存在。证明：如果 $\frac{\partial f}{\partial x}$ 和 $\frac{\partial f}{\partial y}$ 中有一个在 $(x_0, y_0)$ 处连续，那么 $f(x, y)$ 在 $(x_0, y_0)$ 可微。
\end{problem}

\begin{myproof}
    不妨设 $\frac{\partial f}{\partial x}$ 在 $(x_0, y_0)$ 处连续。考虑全增量
    \begin{equation}
        \Delta z = f(x_0 + \Delta x, y_0 + \Delta y) - f(x_0, y_0).
        \label{eq:diff-def-1}
    \end{equation}
    将其拆分为
    \begin{equation}
        \Delta z = [f(x_0 + \Delta x, y_0 + \Delta y) - f(x_0, y_0 + \Delta y)] + [f(x_0, y_0 + \Delta y) - f(x_0, y_0)].
        \label{eq:increment-split}
    \end{equation}
    对 \cref{eq:increment-split} 中第一项关于变量 $x$ 使用一元拉格朗日中值定理，对第二项利用 $f_y(x_0, y_0)$ 的存在性：
    \begin{equation}
        \begin{aligned}
            \Delta z & = f_x(x_0 + \theta \Delta x, y_0 + \Delta y) \Delta x + [f_y(x_0, y_0) \Delta y + o(\Delta y)]    \\
                     & = f_x(x_0 + \theta \Delta x, y_0 + \Delta y) \Delta x + f_y(x_0, y_0) \Delta y + \alpha \Delta y,
        \end{aligned}
        \label{eq:mvt-apply}
    \end{equation}
    其中 $0 < \theta < 1$，且当 $\Delta y \to 0$ 时，$\alpha \to 0$。由于 $f_x$ 在 $(x_0, y_0)$ 处连续，当 $\rho = \sqrt{(\Delta x)^2 + (\Delta y)^2} \to 0$ 时，有
    \begin{equation}
        f_x(x_0 + \theta \Delta x, y_0 + \Delta y) = f_x(x_0, y_0) + \beta,
        \label{eq:fx-cont}
    \end{equation}
    其中 $\beta \to 0$。代入 \cref{eq:mvt-apply} 得
    \begin{equation}
        \Delta z = f_x(x_0, y_0) \Delta x + f_y(x_0, y_0) \Delta y + \beta \Delta x + \alpha \Delta y.
        \label{eq:final-form}
    \end{equation}
    注意到
    \begin{equation}
        \frac{|\beta \Delta x + \alpha \Delta y|}{\sqrt{(\Delta x)^2 + (\Delta y)^2}} \leqslant |\beta| \frac{|\Delta x|}{\rho} + |\alpha| \frac{|\Delta y|}{\rho} \leqslant |\beta| + |\alpha| \to 0.
        \label{eq:limit-zero}
    \end{equation}
    故 $\beta \Delta x + \alpha \Delta y = o(\rho)$。因此 $f(x, y)$ 在 $(x_0, y_0)$ 处可微。
\end{myproof}

\begin{problem}{变量分离型偏微分方程}{Separation-of-Variables-PDE}
    设 $u(x, y)$ 在 $\symbb{R}^2$ 上取正值且有二阶连续偏导数。证明：$u$ 满足方程
    \begin{equation}
        u \frac{\partial^2 u}{\partial x \partial y} = \frac{\partial u}{\partial x} \frac{\partial u}{\partial y}
        \label{eq:pde-origin}
    \end{equation}
    的充分必要条件是存在一元函数 $f$ 和 $g$ 使得 $u(x, y) = f(x)g(y)$。
\end{problem}

\begin{myproof}

    必要性：由 $u(x, y) > 0$，令 $v(x, y) = \ln u(x, y)$。对 $v$ 求偏导得
    \begin{equation}
        \frac{\partial v}{\partial x} = \frac{1}{u} \frac{\partial u}{\partial x},
        \label{eq:vx-pde}
    \end{equation}
    进一步关于 $y$ 求导：
    \begin{equation}
        \frac{\partial^2 v}{\partial x \partial y} = \frac{\partial}{\partial y} \left( \frac{1}{u} \frac{\partial u}{\partial x} \right) = \frac{u \frac{\partial^2 u}{\partial y \partial x} - \frac{\partial u}{\partial y} \frac{\partial u}{\partial x}}{u^2}.
        \label{eq:vxy-pde}
    \end{equation}
    若 \cref{eq:pde-origin} 成立，则有 $\frac{\partial^2 v}{\partial x \partial y} = 0$。积分得
    \begin{equation}
        \frac{\partial v}{\partial x} = A(x) \implies v(x, y) = \int A(x) \d x + \Psi(y) = \Phi(x) + \Psi(y).
        \label{eq:v-integral}
    \end{equation}
    从而有
    \begin{equation}
        u(x, y) = \e^{\Phi(x) + \Psi(y)} = \e^{\Phi(x)} \cdot \e^{\Psi(y)}.
        \label{eq:u-sep-form}
    \end{equation}
    取 $f(x) = \e^{\Phi(x)}$，$g(y) = \e^{\Psi(y)}$ 即证。

    充分性：设 $u(x, y) = f(x)g(y)$。直接计算其偏导数：
    \begin{equation}
        \frac{\partial u}{\partial x} = f'(x)g(y), \quad \frac{\partial u}{\partial y} = f(x)g'(y), \quad \frac{\partial^2 u}{\partial x \partial y} = f'(x)g'(y).
        \label{eq:suff-deriv}
    \end{equation}
    将其代入方程左、右两端：
    \begin{equation}
        \begin{aligned}
            u \frac{\partial^2 u}{\partial x \partial y}                & = [f(x)g(y)] \cdot [f'(x)g'(y)] = f(x)f'(x)g(y)g'(y), \\
            \frac{\partial u}{\partial x} \frac{\partial u}{\partial y} & = [f'(x)g(y)] \cdot [f(x)g'(y)] = f(x)f'(x)g(y)g'(y).
        \end{aligned}
        \label{eq:suff-verify}
    \end{equation}
    两端相等，故方程 \cref{eq:pde-origin} 成立。
\end{myproof}


\begin{problem}{中值定理与向量值函数}{Vector-Mean-Value-Theorem}
    设 $\symbfit{r}(t) = x(t)\symbfit{i} + y(t)\symbfit{j}$，$t \in [a, b]$ 有连续的导数。证明：存在 $\theta \in (a, b)$，使得
    \begin{equation}
        |\symbfit{r}(b) - \symbfit{r}(a)| \le |\symbfit{r}'(\theta)|(b - a). \label{eq:target-inequality}
    \end{equation}
    提示：令 $\varphi(t) = (x(b) - x(a))x(t) + (y(b) - y(a))y(t)$，对函数 $\varphi(t)$ 使用微分中值定理，并利用 Cauchy-Schwarz（柯西-施瓦茨）不等式：$(a_1b_1 + a_2b_2)^2 \le (a_1^2 + a_2^2)(b_1^2 + b_2^2)$。
\end{problem}

\begin{myproof}

    构造辅助函数 $\varphi(t) = (x(b) - x(a))x(t) + (y(b) - y(a))y(t)$。由于 $x(t)$、$y(t)$ 在 $[a, b]$ 上具有连续的导数，故 $\varphi(t)$ 在 $[a, b]$ 上连续，在 $(a, b)$ 内可导。
    根据 Lagrange 中值定理，存在 $\theta \in (a, b)$ 使得
    \begin{equation}
        \varphi(b) - \varphi(a) = \varphi'(\theta)(b - a). \label{eq:lagrange}
    \end{equation}
    计算 $\varphi(b) - \varphi(a)$ 可得
    \begin{equation}
        \begin{aligned}
            \varphi(b) - \varphi(a) & = (x(b) - x(a))x(b) + (y(b) - y(a))y(b) - \left[ (x(b) - x(a))x(a) + (y(b) - y(a))y(a) \right] \\
                                    & = (x(b) - x(a))^2 + (y(b) - y(a))^2                                                            \\
                                    & = |\symbfit{r}(b) - \symbfit{r}(a)|^2.
        \end{aligned} \label{eq:left-side}
    \end{equation}
    计算 $\varphi'(\theta)$ 可得
    \begin{equation}
        \varphi'(\theta) = (x(b) - x(a))x'(\theta) + (y(b) - y(a))y'(\theta). \label{eq:derivative}
    \end{equation}
    由 Cauchy-Schwarz 不等式可知
    \begin{equation}
        \begin{aligned}
            \varphi'(\theta) & \le \sqrt{(x(b) - x(a))^2 + (y(b) - y(a))^2} \cdot \sqrt{(x'(\theta))^2 + (y'(\theta))^2} \\
                             & = |\symbfit{r}(b) - \symbfit{r}(a)| \cdot |\symbfit{r}'(\theta)|.
        \end{aligned} \label{eq:cauchy-schwarz}
    \end{equation}
    将 \cref{eq:left-side} 和 \cref{eq:cauchy-schwarz} 代入 \cref{eq:lagrange}，有
    \begin{equation}
        |\symbfit{r}(b) - \symbfit{r}(a)|^2 \le |\symbfit{r}(b) - \symbfit{r}(a)| \cdot |\symbfit{r}'(\theta)| (b - a).
    \end{equation}
    若 $|\symbfit{r}(b) - \symbfit{r}(a)| = 0$，不等式显然成立；若 $|\symbfit{r}(b) - \symbfit{r}(a)| \ne 0$，两边同除以 $|\symbfit{r}(b) - \symbfit{r}(a)|$ 即得
    \begin{equation}
        \boxed{|\symbfit{r}(b) - \symbfit{r}(a)| \le |\symbfit{r}'(\theta)|(b - a)}
    \end{equation}
\end{myproof}


\begin{problem}{偏导数方程}{Partial-Differential-Equation}
    设 $f(x,y,z)$ 在 $\mathbb{R}^3$ 上有一阶连续偏导数，且满足 $\frac{\partial f}{\partial x} = \frac{\partial f}{\partial y} = \frac{\partial f}{\partial z}$。 如果 $f(x,0,0) > 0$ 对任意的 $x \in \mathbb{R}$ 成立，求证：对任意的 $(x,y,z) \in \mathbb{R}^3$，也有 $f(x,y,z) > 0$。
\end{problem}

\begin{myproof}
    构造辅助函数 $g(t) = f(x+t, y-t, z)$。由复合函数求导法则可得
    \begin{equation}
        \frac{\d g}{\d t} = \frac{\partial f}{\partial x} \cdot 1 + \frac{\partial f}{\partial y} \cdot (-1) = \frac{\partial f}{\partial x} - \frac{\partial f}{\partial y} = 0.
    \end{equation}
    故 $g(t)$ 在 $\mathbb{R}$ 上为常值函数，则 $g(0) = g(y)$，即
    \begin{equation}
        \label{eq:step1}
        f(x,y,z) = f(x+y, 0, z).
    \end{equation}
    同理，构造辅助函数 $h(t) = f(x+y+t, 0, z-t)$，则有
    \begin{equation}
        \frac{\d h}{\d t} = \frac{\partial f}{\partial x} \cdot 1 + \frac{\partial f}{\partial z} \cdot (-1) = \frac{\partial f}{\partial x} - \frac{\partial f}{\partial z} = 0.
    \end{equation}
    故 $h(0) = h(z)$，即
    \begin{equation}
        \label{eq:step2}
        f(x+y, 0, z) = f(x+y+z, 0, 0).
    \end{equation}
    结合 \cref{eq:step1} 与 \cref{eq:step2} 可得，对任意 $(x,y,z) \in \mathbb{R}^3$，均有
    \begin{equation}
        f(x,y,z) = f(x+y+z, 0, 0).
    \end{equation}
    令 $u = x+y+z \in \mathbb{R}$。根据题设条件，对任意 $u \in \mathbb{R}$，均有 $f(u,0,0) > 0$。因此
    \begin{equation}
        f(x,y,z) = f(u, 0, 0) > 0.
    \end{equation}
\end{myproof}


\begin{problem}{多元函数最值}{Multivariable-Extrema}
    求函数 $f(x,y) = x^2 + xy^2 - x$ 在区域 $D = \{(x,y) \mid x^2 + y^2 \leqslant 2\}$ 上的最大值和最小值。
\end{problem}

\begin{solution}
    首先寻找区域 $D$ 内部的驻点。令其偏导数为零：
    \begin{equation}
        \label{eq:stationary-points}
        \left\{ \begin{aligned}
            f_x & = 2x + y^2 - 1 = 0, \\
            f_y & = 2xy = 0.
        \end{aligned} \right.
    \end{equation}
    由 $f_y = 0$ 可得 $x=0$ 或 $y=0$。

    \begin{enumerate}
        \item 若 $x=0$，代入 $f_x=0$ 得 $y^2 = 1$，即 $y = \pm 1$。得到驻点 $(0, 1)$、$(0, -1)$，均在 $D$ 内。此时 $f(0, \pm 1) = 0$。
        \item 若 $y=0$，代入 $f_x=0$ 得 $x = \frac{1}{2}$。得到驻点 $(\frac{1}{2}, 0)$，在 $D$ 内。此时 $f(\frac{1}{2}, 0) = -\frac{1}{4}$。
    \end{enumerate}

    再考察边界 $x^2 + y^2 = 2$。将 $y^2 = 2 - x^2$ 代入原函数，得一元函数
    \begin{equation}
        g(x) = x^2 + x(2 - x^2) - x = -x^3 + x^2 + x, \quad x \in [-\sqrt{2}, \sqrt{2}].
    \end{equation}
    求导得 $g'(x) = -3x^2 + 2x + 1 = 0$，解得 $x_1 = 1$，$x_2 = -\frac{1}{3}$。计算相关点的函数值：
    \begin{enumerate}
        \item 当 $x=1$ 时，$g(1) = 1$。
        \item 当 $x=-\frac{1}{3}$ 时，$g(-\frac{1}{3}) = \frac{1}{27} + \frac{1}{9} - \frac{1}{3} = -\frac{5}{27}$。
        \item 考察边界端点 $x = \pm \sqrt{2}$：
              \begin{equation}
                  g(\sqrt{2}) = 2 - \sqrt{2}, \quad g(-\sqrt{2}) = 2 + \sqrt{2}.
              \end{equation}
    \end{enumerate}

    比较所有候选值：$0$、$-\frac{1}{4}$、$-\frac{5}{27}$、$1$、$2-\sqrt{2}$、$2+\sqrt{2}$。易见最大值为 $2+\sqrt{2}$。由于 $-\frac{1}{4} < -\frac{5}{27}$，故最小值为 $-\frac{1}{4}$。
    \begin{equation}
        \boxed{f_{\max} = 2 + \sqrt{2}, \quad f_{\min} = -\frac{1}{4}}
    \end{equation}
\end{solution}


\begin{problem}{Lagrange 乘数法}{Lagrange-Multipliers-Inequality}
    设 $x_1, x_2, \dots, x_n$ 是正数，且 $x_1 + x_2 + \dots + x_n = n$。用 Lagrange 乘数法证明
    \begin{equation}
        x_1 x_2 \dots x_n \left( \frac{1}{x_1} + \frac{1}{x_2} + \dots + \frac{1}{x_n} \right) \leqslant n,
    \end{equation}
    等号当且仅当 $x_1 = x_2 = \dots = x_n = 1$ 时成立。
\end{problem}

\begin{myproof}
    令 $f(x_1, x_2, \dots, x_n) = \sum_{i=1}^n \frac{\prod_{j=1}^n x_j}{x_i}$，约束条件为 $g(x_1, x_2, \dots, x_n) = \sum_{i=1}^n x_i - n = 0$。

    构造 Lagrange 函数：
    \begin{equation}\label{eq:Lagrange-Function}
        L(x_1, x_2, \dots, x_n, \lambda) = f(x_1, x_2, \dots, x_n) + \lambda \left( \sum_{i=1}^n x_i - n \right).
    \end{equation}

    令 $P = \prod_{j=1}^n x_j$，则 $f = \sum_{i=1}^n \frac{P}{x_i}$。对各变量求偏导数：
    \begin{equation}\label{eq:Partial-Deriv}
        \frac{\partial L}{\partial x_k} = \frac{\partial f}{\partial x_k} + \lambda = \sum_{i \neq k} \frac{P}{x_i x_k} + \lambda = \frac{1}{x_k} \left( f - \frac{P}{x_k} \right) + \lambda = 0, \quad k = 1, 2, \dots, n.
    \end{equation}

    整理 \cref{eq:Partial-Deriv} 可得：
    \begin{equation}\label{eq:Condition-1}
        f - \frac{P}{x_k} + \lambda x_k = 0.
    \end{equation}

    对 \cref{eq:Condition-1} 关于 $k$ 求和：
    \begin{equation}
        \sum_{k=1}^n \left( f - \frac{P}{x_k} + \lambda x_k \right) = nf - f + \lambda n = (n-1)f + \lambda n = 0 \Rightarrow \lambda = -\frac{n-1}{n} f.
    \end{equation}

    代入 \cref{eq:Condition-1} 得：
    \begin{equation}
        f - \frac{P}{x_k} - \frac{n-1}{n} f x_k = 0 \Rightarrow \frac{n-1}{n} f x_k^2 - f x_k + P = 0.
    \end{equation}

    由于该关于 $x_k$ 的二次方程至多有两个正根，结合对称性及约束条件 $\sum x_i = n$，可知驻点为 $x_1 = x_2 = \dots = x_n = 1$。此时：
    \begin{equation}
        f(1, 1, \dots, 1) = \sum_{i=1}^n 1 = n.
    \end{equation}

    考虑边界情况，若某些 $x_i \to 0$，则 $P \to 0$，从而 $f$ 的极限值不大于 $n$。故 $f$ 在该点取得最大值。

    因此，原不等式成立：
    \begin{equation}
        \boxed{x_1 x_2 \dots x_n \left( \frac{1}{x_1} + \frac{1}{x_2} + \dots + \frac{1}{x_n} \right) \leqslant n}
    \end{equation}
\end{myproof}



\begin{problem}{Jensen 不等式}{Power-Mean-Inequality}
    设 $a_i \ge 0$ ($i = 1, 2, \dots, n$)、$p > 1$。证明：
    \begin{equation}
        \frac{a_1 + \cdots + a_n}{n} \le \left( \frac{a_1^p + \cdots + a_n^p}{n} \right)^{1/p},
    \end{equation}
    并讨论等号成立的条件。
\end{problem}

\begin{myproof}
    考虑函数 $f(x) = x^p$，$x \in [0, +\infty)$。由于 $p > 1$，计算其二阶导数得
    \begin{equation}
        f''(x) = p(p-1)x^{p-2}. \label{eq:second-derivative}
    \end{equation}
    当 $x > 0$ 时，$f''(x) > 0$，故 $f(x)$ 在 $[0, +\infty)$ 上为严格凸函数。根据 Jensen 不等式，有
    \begin{equation}
        f\left( \frac{1}{n} \sum_{i=1}^{n} a_i \right) \le \frac{1}{n} \sum_{i=1}^{n} f(a_i), \label{eq:jensen-apply}
    \end{equation}
    代入 $f(x)$ 的表达式得
    \begin{equation}
        \left( \frac{a_1 + \cdots + a_n}{n} \right)^p \le \frac{a_1^p + \cdots + a_n^p}{n}.
    \end{equation}
    由于 $p > 1$，对上式两端同时取 $1/p$ 次幂，不等号方向不变，即
    \begin{equation}
        \frac{a_1 + \cdots + a_n}{n} \le \left( \frac{a_1^p + \cdots + a_n^p}{n} \right)^{1/p}.
    \end{equation}
    由于 $f(x)$ 为严格凸函数，根据 \cref{eq:jensen-apply}，等号成立的充分必要条件为
    \begin{equation}
        \boxed{a_1 = a_2 = \cdots = a_n.}
    \end{equation}
\end{myproof}


\begin{problem}{外微分与雅可比行列式}{Differential-Form-Jacobian}
    设 $y_i = y_i(x_1, \dots, x_n)$ ($i = 1, 2, \dots, n$) 是 $n$ 个可微的 $n$ 元函数，证明：
    \begin{equation}
        \d y_1 \wedge \d y_2 \wedge \dots \wedge \d y_n = \frac{\partial(y_1, y_2, \dots, y_n)}{\partial(x_1, x_2, \dots, x_n)} \d x_1 \wedge \d x_2 \wedge \dots \wedge \d x_n.
    \end{equation}
\end{problem}

\begin{myproof}
    根据全微分的定义，对于每一个 $i = 1, 2, \dots, n$，有
    \begin{equation}
        \d y_i = \sum_{j=1}^n \frac{\partial y_i}{\partial x_j} \d x_j. \label{eq:diff-definition-1}
    \end{equation}

    将 \cref{eq:diff-definition-1} 代入左端的外积表达式中，利用外积的线性性质，得
    \begin{equation}
        \begin{aligned}
            \d y_1 \wedge \d y_2 \wedge \dots \wedge \d y_n & = \left( \sum_{j_1=1}^n \frac{\partial y_1}{\partial x_{j_1}} \d x_{j_1} \right) \wedge \left( \sum_{j_2=1}^n \frac{\partial y_2}{\partial x_{j_2}} \d x_{j_2} \right) \wedge \dots \wedge \left( \sum_{j_n=1}^n \frac{\partial y_n}{\partial x_{j_n}} \d x_{j_n} \right) \\
                                                            & = \sum_{j_1, j_2, \dots, j_n = 1}^n \frac{\partial y_1}{\partial x_{j_1}} \frac{\partial y_2}{\partial x_{j_2}} \dots \frac{\partial y_n}{\partial x_{j_n}} \d x_{j_1} \wedge \d x_{j_2} \wedge \dots \wedge \d x_{j_n}.
        \end{aligned}
    \end{equation}

    利用外积的反对称性，当指标序列 $(j_1, j_2, \dots, j_n)$ 中存在重复数字时，对应的微分形式乘积为零。因此，求和仅需遍历 $\{1, 2, \dots, n\}$ 的所有排列 $\sigma$。令 $j_i = \sigma(i)$，则有
    \begin{equation}
        \d x_{\sigma(1)} \wedge \d x_{\sigma(2)} \wedge \dots \wedge \d x_{\sigma(n)} = \operatorname{sgn}(\sigma) \d x_1 \wedge \d x_2 \wedge \dots \wedge \d x_n.
    \end{equation}

    由此可得
    \begin{equation}
        \begin{aligned}
            \d y_1 \wedge \d y_2 \wedge \dots \wedge \d y_n & = \left( \sum_{\sigma \in S_n} \operatorname{sgn}(\sigma) \frac{\partial y_1}{\partial x_{\sigma(1)}} \frac{\partial y_2}{\partial x_{\sigma(2)}} \dots \frac{\partial y_n}{\partial x_{\sigma(n)}} \right) \d x_1 \wedge \d x_2 \wedge \dots \wedge \d x_n \\
                                                            & = \det \left( \frac{\partial y_i}{\partial x_j} \right) \d x_1 \wedge \d x_2 \wedge \dots \wedge \d x_n.
        \end{aligned}
    \end{equation}

    根据雅可比行列式的定义
    \begin{equation}
        \frac{\partial(y_1, y_2, \dots, y_n)}{\partial(x_1, x_2, \dots, x_n)} = \det \left( \frac{\partial y_i}{\partial x_j} \right).
    \end{equation}

    证毕，即得
    \begin{equation}
        \boxed{\d y_1 \wedge \d y_2 \wedge \dots \wedge \d y_n = \frac{\partial(y_1, y_2, \dots, y_n)}{\partial(x_1, x_2, \dots, x_n)} \d x_1 \wedge \d x_2 \wedge \dots \wedge \d x_n}
    \end{equation}
\end{myproof}
